% GENERATED FILE. Edit canonical lessons, then run npm run generate:books.
% canonical-source-hash: fnv1a64:f9dd484d10954310
% canonical-lessons: LA-C211-omnino, LA-C211-vix, LA-C211-sponte, LA-C211-frustra, LA-C211-rursus

\chapter{Altogether, Hardly, Of one's own accord, In vain, Back again}
\label{ch:la-altogether-hardly-of-one-s-own-accord-in-vain-back-again}
\hlchaptermodality{211}

\begin{chapteropening}
\textbf{By the end of this chapter:} \emph{I can say altogether, hardly, of one's own accord, in vain and back again.}

Complete the last lesson of chapter 211: I can say altogether, hardly, of one's own accord, in vain and back again.
\end{chapteropening}

\section[omnīnō]{omnīnō — altogether}

\label{lesson:LA-C211-omnino}



\begin{quote}
Before the new one: say the Latin for about, then the Latin for only; so much.
\end{quote}

\subsection*{You'll want to know: omnīnō}
\textbf{omnīnō} — \textquotedblleft{}altogether\textquotedblright{}.

\begin{cousinweb}
\textbf{omnīnō} is built on \emph{omnis}, \textquotedblleft{}all\textquotedblright{} (Lewis \& Short): altogether. Where \emph{omnis} comes from is uncertain.
\end{cousinweb}

\begin{grammarlens}[title={What you've built}]
One more word for this chapter.
\end{grammarlens}

\subsection*{Your turn}
\begin{itemize}
\raggedright
  \item \emph{Say it:} \emph{omnīnō}
  \item \emph{Say it:} \emph{omnīnō}, once more
  \item \emph{Say it:} say \emph{omnīnō}
  \item \emph{Recall:} say the Latin for about, then the Latin for only; so much, then say \emph{omnīnō} again
\end{itemize}

\begin{tcolorbox}[breakable,colback=teal!4,colframe=teal!35!black,title={Before you move on}]
What does omnīnō mean? (\textquotedblleft{}Altogether\textquotedblright{}.) Say it once more.
\end{tcolorbox}
\section[vix]{vix — hardly}

\label{lesson:LA-C211-vix}



\begin{quote}
Before the new one: say the Latin for only; so much, then the Latin for altogether.
\end{quote}

\subsection*{You'll want to know: vix}
\textbf{vix} — \textquotedblleft{}hardly\textquotedblright{}.

\begin{cousinweb}
\textbf{vix} has no sure origin; Lewis \& Short suggest the root of \emph{vincō}, \textquotedblleft{}to win\textquotedblright{}: won with effort. That is uncertain.
\end{cousinweb}

\begin{grammarlens}[title={What you've built}]
One more word for this chapter.
\end{grammarlens}

\subsection*{Your turn}
\begin{itemize}
\raggedright
  \item \emph{Say it:} \emph{vix}
  \item \emph{Say it:} \emph{vix}, once more
  \item \emph{Say it:} say \emph{vix}
  \item \emph{Recall:} say the Latin for only; so much, then the Latin for altogether, then say \emph{vix} again
\end{itemize}

\begin{tcolorbox}[breakable,colback=teal!4,colframe=teal!35!black,title={Before you move on}]
What does vix mean? (\textquotedblleft{}Hardly\textquotedblright{}.) Say it once more.
\end{tcolorbox}
\section[sponte]{sponte — of one's own accord}

\label{lesson:LA-C211-sponte}



\begin{quote}
Before the new one: say the Latin for altogether, then the Latin for hardly.
\end{quote}

\subsection*{You'll want to know: sponte}
\textbf{sponte} — \textquotedblleft{}of one's own accord\textquotedblright{}.

\begin{cousinweb}
\textbf{sponte} is the ablative of an old noun for \textquotedblleft{}free will\textquotedblright{}; Lewis \& Short link it with \emph{spondeō}, \textquotedblleft{}to pledge\textquotedblright{}, but that is uncertain.
\end{cousinweb}

\begin{grammarlens}[title={What you've built}]
One more word for this chapter.
\end{grammarlens}

\subsection*{Your turn}
\begin{itemize}
\raggedright
  \item \emph{Say it:} \emph{sponte}
  \item \emph{Say it:} \emph{sponte}, once more
  \item \emph{Say it:} say \emph{sponte}
  \item \emph{Recall:} say the Latin for altogether, then the Latin for hardly, then say \emph{sponte} again
\end{itemize}

\begin{tcolorbox}[breakable,colback=teal!4,colframe=teal!35!black,title={Before you move on}]
What does sponte mean? (\textquotedblleft{}Of one's own accord\textquotedblright{}.) Say it once more.
\end{tcolorbox}
\section[frūstrā]{frūstrā — in vain}

\label{lesson:LA-C211-frustra}



\begin{quote}
Before the new one: say the Latin for hardly, then the Latin for of one's own accord.
\end{quote}

\subsection*{You'll want to know: frūstrā}
\textbf{frūstrā} — \textquotedblleft{}in vain\textquotedblright{}.

\begin{cousinweb}
\textbf{frūstrā} is usually linked with \emph{fraus}, \textquotedblleft{}deceit, harm\textquotedblright{} (Lewis \& Short): to one's harm, so in vain.
\end{cousinweb}

\begin{grammarlens}[title={What you've built}]
One more word for this chapter.
\end{grammarlens}

\subsection*{Your turn}
\begin{itemize}
\raggedright
  \item \emph{Say it:} \emph{frūstrā}
  \item \emph{Say it:} \emph{frūstrā}, once more
  \item \emph{Say it:} say \emph{frūstrā}
  \item \emph{Recall:} say the Latin for hardly, then the Latin for of one's own accord, then say \emph{frūstrā} again
\end{itemize}

\begin{tcolorbox}[breakable,colback=teal!4,colframe=teal!35!black,title={Before you move on}]
What does frūstrā mean? (\textquotedblleft{}In vain\textquotedblright{}.) Say it once more.
\end{tcolorbox}
\section[rūrsus]{rūrsus — back again, once more}

\label{lesson:LA-C211-rursus}



\begin{quote}
Before the new one: say the Latin for of one's own accord, then the Latin for in vain.
\end{quote}

\subsection*{You'll want to know: rūrsus}
\textbf{rūrsus} — \textquotedblleft{}back again, once more\textquotedblright{}.

\begin{cousinweb}
\textbf{rūrsus} is contracted from \emph{revorsus}, \textquotedblleft{}turned back\textquotedblright{} (Lewis \& Short), the participle of \emph{revertō}, \textquotedblleft{}to turn back\textquotedblright{}.
\end{cousinweb}

\begin{grammarlens}[title={What you've built}]
One more word for this chapter.
\end{grammarlens}

\subsection*{Your turn}
\begin{itemize}
\raggedright
  \item \emph{Say it:} \emph{rūrsus}
  \item \emph{Say it:} \emph{rūrsus}, once more
  \item \emph{Say it:} say \emph{rūrsus}
  \item \emph{Recall:} say the Latin for of one's own accord, then the Latin for in vain, then say \emph{rūrsus} again
\end{itemize}

\begin{tcolorbox}[breakable,colback=teal!4,colframe=teal!35!black,title={Before you move on}]
What does rūrsus mean? (\textquotedblleft{}Back again, once more\textquotedblright{}.) Say it once more.
\end{tcolorbox}
